\chapter{APPENDIX: TEXT PROCESSING AND LAYOUT IN A THESIS}\label{appendix:text_processing_and_layout_in_a_thesis}
Good text processing skills make writing your final thesis and using this template easier. Therefore, you should make sure you have the sufficient basic skills to edit long documents with text processing software before you start. This involves applying the styles, understanding automatic referencing and knowing how to divide your text into sections.

The essential thing is to understand the following basics of Word:
% Create bullet point list. Noitemsep option decreases space between items.
\begin{itemize}[noitemsep]
	\item Do not modify your layout by adding consecutive spaces or line breaks. If you need to press Enter or the space bar more than once, you are probably doing something wrong. When you want to start a new paragraph, press Enter once at the end of the sentence and use styles to create a space between the paragraphs.
	\item Do not do numbering manually. Word has efficient automatised tools for this. Section numbering and page numbering are applied in this template but also automatised numbering of figures/tables is possible. They keep the numbers in the right order even if you modify, add or remove information.
	\item Do not add hyphenation at the end of a line manually. Word’s automatic hyphenation tool can be used in this template. If you need to add more hyphens, select manual hyphenation in Word. The automatic hyphenation is usually turned off, but you can activate it yourself.
\end{itemize}

\underline{Line spacing, font, margins, alignment, page numbering and headings}

Official layout guidelines state that the line spacing should be 1.5 except for the abstract and possible direct citations, where the spacing is 1. You can choose from two fonts: Times New Roman (12 pt) or Arial (11 pt). This template has been written in Times New Roman 12 pt (Style LUT Normal). Leave an empty line before and after tables and figures.

Leave the following margins:
\begin{itemize}[noitemsep]
	\item top and left 35 mm
	\item bottom and right 20 mm. 
\end{itemize}

The page number of the title page is 1, but the page numbering should not be visible before the first page of the table of contents. Place the page numbers at the top of the page, either centred or in the right-hand corner. Page numbering ends on the final page of the reference list: appendices do not have page numbers unless the appendix is multiple pages.

% In LaTeX percentage sign must be prefixed with backslash.
Always use heading styles in your headings (LUT Heading 1, LUT Heading 2, LUT Heading 3). Always place chapter headings (LUT Heading 1) on a fresh page. If you add large figures or tables, remember to check that the empty space after the figure or image covers no more than 20\% of the page.

Always use heading styles in your headings (LUT Heading 1, LUT Heading 2, LUT Heading 3). Always place chapter headings (LUT Heading 1) on a fresh page. If you add large figures or tables, remember to check that the empty space after the figure or image covers no more than 20\% of the page.

The thesis should include no more than three heading levels, and the headings should progress in a logical order (first LUT Heading 1, then LUT Heading 2, etc.). If you need even more detailed subheadings, do not number them, and leave them out of the table of contents. Concise headings that describe the text sufficiently are the best. You can use question marks or exclamation marks, but do not add a period if the heading is a regular sentence.

A heading cannot be followed by a heading. Always write something between them. For instance, there must be text between headings 1 and 1.1.

Lists are a good way to express things clearly. Use the same type of bullet or symbol in lists throughout your thesis. A section should never end in a list. There should always be two or three sentences after a list.